\section{O CRESCIMENTO DA POPULAÇÃO DE ANIMAIS DOMÉSTICOS NO BRASIL}

O cenário demográfico brasileiro tem registrado uma expansão contínua no número de animais de estimação, redefinindo as relações afetivas e a rotina das famílias atuais. Segundo levantamentos da Pesquisa Nacional por Amostra de Domicílios Contínua realizada pelo Instituto Brasileiro de Geografia e Estatística (IBGE, 2023), quase metade dos lares brasileiros conta com pelo menos um animal integrado ao ambiente familiar. Em números absolutos, dados consolidados pelo Instituto Pet Brasil (2023) apontam que a população de animais de estimação no país ultrapassou a expressiva marca de 149 milhões de animais. Essa convivência diária consolida os animais como verdadeiros membros do convívio familiar, tornando a guarda responsável e o bem-estar animal assuntos de grande relevância para os centros urbanos (INSTITUTO PET BRASIL, 2023).

No entanto, à medida que a população de animais cresce nos centros urbanos, aumenta também a ocorrência de fugas e desaparecimentos. Quando um animal de estimação escapa para a rua, as primeiras horas após o sumiço são determinantes para o sucesso do resgate. Desorientado e assustado com o barulho do trânsito, o animal tende a se afastar rapidamente de casa, percorrendo longas distâncias sem rumo e enfrentando perigos constantes, como atropelamentos, falta de alimento, chuvas fortes e agressões físicas. Essa vulnerabilidade atinge níveis ainda mais graves durante situações de calamidade pública e desastres naturais, a exemplo das enchentes históricas que afetaram o estado do Rio Grande do Sul em 2024, quando milhares de animais foram separados abruptamente de seus tutores, sobrecarregando voluntários e protetores independentes (CRMV-RS, 2024; UOL, 2024).

Apesar da gravidade desse problema, as formas mais comuns utilizadas pelas pessoas para encontrar animais perdidos (principalmente colar cartazes em postes e publicar em grupos de redes sociais) apresentam sérias limitações práticas. A colagem manual de cartazes de papel em postes, muros e estabelecimentos comerciais exige gastos com impressão, demanda tempo excessivo para ser feita e possui alcance restrito a poucas quadras. Além disso, os cartazes impressos estragam-se rapidamente com a chuva e o vento, perdendo a legibilidade em poucos dias.

Com a difusão da internet, os tutores passaram a utilizar redes sociais generalistas, como Facebook e Instagram, e grupos de mensagens instantâneas no WhatsApp. Embora alcancem um número maior de pessoas, esses canais não foram projetados para situações de emergência e resgate. Neles, as postagens perdem alcance aceleradamente à medida que novas mensagens entram no fluxo de postagens, as informações ficam espalhadas em comentários paralelos e não existem recursos de busca organizados por mapa ou proximidade geográfica. Outro problema crítico envolve a privacidade: para viabilizar o contato, os tutores sentem-se forçados a publicar seus números de telefone e endereços pessoais de forma aberta, ficando expostos a trotes, tentativas de extorsão e fraudes financeiras.

Por outro lado, o avanço tecnológico estabeleceu o smartphone como o principal meio de acesso à internet nos lares brasileiros (CETIC.br, 2023; CNDL; SPC BRASIL, 2024). A presença constante dos aparelhos celulares permite conectar moradores em redes de apoio comunitário em tempo real (CASTELLS, 2022). No âmbito da computação, pesquisas indicam que aplicações móveis dedicadas com recursos de geolocalização e mapas interativos apresentam alto potencial para agilizar a localização de animais perdidos (SOUZA et al., 2022). Ao permitir que as pessoas visualizem publicações por proximidade, registrem fotos com o ponto exato da publicação e conversem por canais seguros sem divulgar dados particulares, a tecnologia móvel supera as falhas dos cartazes de papel e das publicações dispersas em redes sociais.

Para construir uma solução tecnológica que responda com precisão a essas necessidades, foi fundamental consultar diretamente os tutores e analisar as plataformas disponíveis no mercado. O Capítulo 3 apresenta a metodologia adotada na pesquisa, detalhando o levantamento de campo realizado com a comunidade e o estudo comparativo de sistemas similares.
\clearpage
